\chapter{Understanding the Solver output}
\label{chap:solverout}

This Appendix describes, in more detail than the text in Chapter \ref{chap:UnderstandingSolver} ``Understanding the Solver,'' how to interpret parts of the output file of the solver. An example from a real, non-trivial run of the solver is presented and examined.

Consider the section of the solver output called ``Final Solution'' near the bottom of the file. This begins with a table of all the measurement data including measurements and residuals. The table is followed by a list of the three entries in the table with the largest standard residuals. Then comes the summary of statistics that is described in Chapter \ref{chap:UnderstandingSolver}.

[Note that the printed precision of numbers in this table is controlled by the solver options \verb+--linprec+ \verb+--angprecM+ and \verb+--angprecP+. In addition, the solver will also print this table at each iteration if the \verb+--verbose+ option is found (also with standard residuals if \verb+--statsAll+). See the solver command line reference Chapter \ref{chap:datref}.]

\newpage
\begin{landscape}
\fontsize{7.2}{7.2}
\begin{Snippet}
 Data and residuals (Final)
 Label                    Meas(m|DMS) Nomin(m|DMS) Pre-R(m|soa) Post-R(m|soa)   Rel-Res   Std-Res  Redund Int-Rel(m|SOA) Ext-Rel-E(m) Ext-Rel-N(m) Ext-Rel-U(m) Ext-Rel-Mag
 Del.T8(VNDP-OSS_TP1)X     -682.22620   -682.22080     -0.00181      -0.00540  -1.47760  -1.60865  0.5659        0.02058      0.01111      0.00349     -0.00251     0.01191
 Del.T8(VNDP-OSS_TP1)Y     1183.31220   1183.31748     -0.00145      -0.00528  -0.45181  -0.53724  0.4760        0.02615     -0.00888      0.00505     -0.00608     0.01189
 Del.T8(VNDP-OSS_TP1)Z     1215.35160   1215.35195     -0.00275      -0.00035  -0.73950  -0.90598  0.4427        0.01981      0.00122      0.00952      0.00449     0.01060
 Del.T9(5N-OSS_TP1)X      -3333.41390  -3333.42650      0.01619       0.01260   1.90033   1.66670  0.8686        0.03005      0.00313      0.00138     -0.00216     0.00404
 Del.T9(5N-OSS_TP1)Y      -3034.89880  -3034.91172      0.01675       0.01292   1.11493   0.95988  0.8963        0.03529     -0.00167      0.00195     -0.00258     0.00364
 Del.T9(5N-OSS_TP1)Z      -6138.44080  -6138.43635     -0.00685      -0.00445  -0.11113  -0.09554  0.9017        0.03134      0.00006      0.00269      0.00135     0.00301
 Del.T10(VNDP-OSS_TP2)X    -639.20240   -639.19906      0.17191      -0.00334  -0.83647  -0.88349  0.5973        0.02182     -0.00004      0.00472     -0.00173     0.01336
 Del.T10(VNDP-OSS_TP2)Y    1143.13750   1143.14303     -2.17457      -0.00553  -0.89714  -1.18186  0.3823        0.03075     -0.00190     -0.00173     -0.00529     0.01603
 Del.T10(VNDP-OSS_TP2)Z    1195.07860   1195.07789     -1.72257       0.00071  -0.22129  -0.29052  0.3888        0.02230     -0.00403      0.00244      0.00299     0.01318
 Del.T11(5N-OSS_TP2)X     -3290.38760  -3290.40476      0.19241       0.01716   2.39202   2.08804  0.8709        0.03236      0.00005      0.00148     -0.00159     0.00419
 Del.T11(5N-OSS_TP2)Y     -3075.06890  -3075.08617     -2.15177       0.01727   1.22928   1.05730  0.9039        0.03847     -0.00034     -0.00002     -0.00195     0.00375
 Del.T11(5N-OSS_TP2)Z     -6158.71370  -6158.71041     -1.72657      -0.00329   0.15504   0.13371  0.8955        0.03227     -0.00103      0.00047      0.00111     0.00329
 Del.T12(VNDP-OSS_TP3)X    -691.84980   -691.84291     -0.37619      -0.00689  -1.57596  -1.56023  0.6833        0.02238      0.00409     -0.00380     -0.00143     0.01174
 Del.T12(VNDP-OSS_TP3)Y    1158.97540   1158.98088     -2.77565      -0.00548   0.01523   0.01859  0.4480        0.03091     -0.00149      0.00792     -0.00441     0.01408
 Del.T12(VNDP-OSS_TP3)Z    1173.01550   1173.01021      0.04978       0.00529   0.88973   1.07754  0.4563        0.02135      0.00279      0.00257      0.00206     0.01116
 Del.T13(5N-OSS_TP3)X     -3343.03330  -3343.04861     -0.35399       0.01531   1.88901   1.61984  0.9083        0.03591      0.00121     -0.00072     -0.00123     0.00354
 Del.T13(5N-OSS_TP3)Y     -3059.23540  -3059.24832     -2.75725       0.01292   0.53935   0.45882  0.9245        0.04236     -0.00015      0.00161     -0.00165     0.00317
 Del.T13(5N-OSS_TP3)Z     -6180.78090  -6180.77809      0.04168      -0.00281   0.15331   0.13104  0.9171        0.03315      0.00061      0.00048      0.00100     0.00286
 Dis.T7(newPOSG-OSS_TP2)     26.00000     26.00413     -1.72816      -0.00413  -0.01376  -0.01124  0.9993        1.26625     -0.00002      0.00002     -0.00001     0.00007
 Han.T14(OSS_TP2-OSS_TP1)  62 41 38.0   62 41 39.5      -9866.8          -1.5      -0.2      -0.3  0.3036       55.19650      0.00262     -0.00296      0.00008     0.00395
 Han.T15(OSS_TP1-OSS_TP2) 311 46 17.9  311 46 20.7      -2440.1          -2.8      -0.4      -0.6  0.3112       51.56637     -0.00120     -0.00346      0.00002     0.00466
 Han.T16(OSS_TP1-OSS_TP3)  69  4 34.7   69  4 41.2       7421.3          -6.5      -0.5      -0.5  0.6983       66.00486     -0.00138     -0.00014     -0.00002     0.00195
 Zan.T17(OSS_TP1-OSS_TP2)  91  0  8.3   90 59 59.7      -1739.2           8.7       0.9       0.9  0.6543       51.56637     -0.00034     -0.00010      0.00243     0.00309
 Zan.T18(OSS_TP1-OSS_TP3)  93  9 35.9   93  9 35.6      -9234.2           0.4       0.0       0.1  0.3217       70.13019      0.00005     -0.00065      0.00403     0.00648
 Zan.T19(OSS_TP2-OSS_TP1)  89  0  8.7   89  0  2.4       1754.2           6.3       0.6       0.6  0.6543       51.56637      0.00034      0.00010     -0.00243     0.00309
 Zan.T20(OSS_TP2-OSS_TP3)  91 36 16.5   91 36 12.9      -5737.0           3.6       0.4       0.4  0.4995       59.81692      0.00047     -0.00024      0.00001     0.00490

 Largest 3 external reliability vector magnitudes:
 Del.T10(VNDP-OSS_TP2)Y    1143.13750   1143.14303     -2.17457      -0.00553  -0.89714  -1.18186  0.3823        0.03075     -0.00190     -0.00173     -0.00529     0.01603
 Del.T12(VNDP-OSS_TP3)Y    1158.97540   1158.98088     -2.77565      -0.00548   0.01523   0.01859  0.4480        0.03091     -0.00149      0.00792     -0.00441     0.01408
 Del.T10(VNDP-OSS_TP2)X    -639.20240   -639.19906      0.17191      -0.00334  -0.83647  -0.88349  0.5973        0.02182     -0.00004      0.00472     -0.00173     0.01336

 Largest 3 standard residuals (** denotes possible blunder: threshold based on Tau distribution and redundancy for each obs)
 Del.T11(5N-OSS_TP2)X     -3290.38760  -3290.40476      0.19241       0.01716   2.39202   2.08804  0.8709        0.03236      0.00005      0.00148     -0.00159     0.00419
 Del.T9(5N-OSS_TP1)X      -3333.41390  -3333.42650      0.01619       0.01260   1.90033   1.66670  0.8686        0.03005      0.00313      0.00138     -0.00216     0.00404
 Del.T13(5N-OSS_TP3)X     -3343.03330  -3343.04861     -0.35399       0.01531   1.88901   1.61984  0.9083        0.03591      0.00121     -0.00072     -0.00123     0.00354

 RMS post-fit relative residual (Final) = 9.37613e-01
 Degrees of freedom = 17; Std dev of unit weight = 1.22461; APV = 1.49967 (Final)
 RMS post-fit raw residual (angles) (Final) = 2.5e-05 rad = 5.1e+00 soa.
 RMS post-fit raw residual (length) (Final) = 9.11888e-03 m.
 Chi-squared test lower bound (0.050): 0.510 < 1.500 pass
 Chi-squared test upper bound (0.950): 1.500 < 1.623 pass
\end{Snippet}
\normalsize
\end{landscape}

\normalsize
Consider each of the columns of this table in turn (some of the column labels have been shortened in the table above to save space). 

\textbf{Column 1 Label.} This is the label that the solver uses to identify the measurements. It is of the form type.label(From-To)[XYZ], which consists of (a) the 3-letter measurement type from the input file (Dis = distance, Zan = zenith angle, etc.), (b) a unique label, which will be a number unless ``tag=label'' appears with the measurement in the input file, and (c) the labels for the Points involved in the measurement separated by dashes, usually just From-To, although horizontal angles will have From-At-To. Delta measurements (always in groups of 3) will also have a coordinate (X,Y,Z) at the end, since the measurement is a 3-vector.

For example, Del.T6(VNDP-TP1)Y is the Y-component of the delta measurement from Point VNDP to Point TP1; Zan.T16(TP1-TP3) is the zenith angle from Point TP1 to Point TP3.

\textbf{Column 2 Meas(m$\mid$DMS).} This is the measurement as read from the input file. Angles are expressed in DMS and linear measurements in meters.

\textbf{Column 3 Nomin(m$\mid$DMS).} This is the nominal or ``predicted'' value of the measurement, meaning the value computed using the current positions of all the Points.

\textbf{Column 4 Pre-R(m$\mid$soa).} The pre-residual is simply the difference ``Measured minus Nominal,'' using the current value for the Nominal (which is shown in column 3). In the final solution (after the last iteration) this means the Nominal as computed from the final adjusted Points. It is called ``pre'' because it is the data computed before (and passed into) the linear least squares algorithm.

\textbf{Column 5 Post-R(m$\mid$soa).} The post-residual is simply the difference ``Measured minus Nominal,'' using the \emph{a priori} value for the Nominal (not shown in the table); it is the total net residual after (hence ``post'') the adjustment. Thus if the \emph{a priori} position of a Point is not very good (e.g. it had to be computed by the solver), this post-residual may very well be large. 

\textbf{Column 6 Rel-Resid.} The relative residual is the output of the least squares solution algorithm at the latest iteration. Thus it indicates how much the measurement is ``in error'' according to the algorithm. This quantity is dimensionless, because the data given to the algorithm (and thus also its output) has been weighted (multiplied) by the inverse measurement covariance matrix. One can think of it roughly as the measurement residual in physical units divided by the given uncertainty in that measurement. (Of course this is a rough characterization because the division is actually a full matrix operation.)

\textbf{Column 7 Std-Resid.} The standard residual, actually a statistical test metric, is defined in appendix \ref{appendices:RedStdResComp}. This test metric is compared to a threshold computed using a 95\% confidence level. Standard residual values exceeding their respective thresholds are marked with ``**'' for the user, indicating a possible blunder within the survey adjustment network (so not necessarily within that measurement). The above table does not have any standard residuals exceeding their respective thresholds.

The three largest magnitude standard residual measurements are repeated in a subsequent table, prior to the statistical summary. Note that the standard residual is undefined for measurements with zero redundancy (which is very undesirable for any measurement in any survey), resulting in two dashes for the standard residual column value. 

\textbf{Column 8 Redund.} Redundancy is a number between 0 and 1 that is a measure of the immunity of the estimation to errors in the subject measurement. This column is the portion of the redundancy of the problem that belongs to this measurement; the sum of all the redundancies is equal to the degrees of freedom of the problem. Thus a large degree of freedom implies geometric strength in the estimation, and the geometry of the network and placement of the measurements are combined to ``divide up'' the total redundancy among the measurements.

It helps to think of redundancy in the two extreme cases. First, consider the case of zero redundancy. Suppose a DXYZ (GNSS vector) involves one fixed Point A and another unknown Point B; then if \emph{no} other measurements involve B, then the adjustment determines B simply by adding the DXYZ to A, and there is no more information pertaining to, and so no freedom to adjust, the measurement to Point B. The redundancy of the DXYZ will be zero (so will the pre-residual and relative residual, the standard residual is undefined).

Second, consider the case of high redundancy. Suppose one measurement involves two Points that also are involved in many other measurements. If you imagine perturbing that measurement, then the algorithm would respond by adjusting many other Points and measurements. The measurement is ``tied'' to many others through a strong network design. Conversely, all the other measurements to which it is tied will serve to improve the determination of the measurement's residual; redundancy in the network is necessary to identify the measurement as discrepant. Ghilani \cite{ghilani} has more discussion, including the derivation of redundancy.

\textbf{Column 9 Int-Rel.} Internal reliability, also known as the minimum detectable bias, is a quantity derived from redundancy. This is the minimum error value for a particular measurement which can be detected on the basis of elevated standard residuals.

\textbf{Column 10-12 Ext-Rel-*(m).} External reliability is calculated by mapping the internal reliability to state-space. Columns 10-12 represent the three components (East-North-Up) of the external reliability vector. 
 
\textbf{Column 13 Ext-Rel-Mag(m).} External reliability magnitude is the magnitude of the measurement's external reliability vector.  It is the Euclidean norm of the East, North, and Up components listed in the preceding columns.

%\section*{Record Types}
%The GeoLab user documentation\cite{geolab} lists the many record types supported by that application.  Many of those record types are not applicable in \textsf{salsa} or are otherwise not supported.  When a GeoLab \verb+.iob+ file is imported in \textsf{salsa}, the command-line application \textsf{iob2lsa} reads the \verb+.iob+ file(s) and generates corresponding \verb+.lsa+ files.  In doing so, the following record types are supported.  GeoLab record types not shown in this table are not supported in \textsf{iob2lsa}; warnings for unsupported record types will be shown in \textsf{salsa}'s status window.
